\chapter{KESIMPULAN DAN SARAN}
\label{chap:kesimpulan-saran}

Bab ini merupakan bagian penutup yang merangkum hasil akhir dari keseluruhan proses penelitian yang telah dijabarkan dan dianalisis pada bab-bab sebelumnya. Pada bagian pertama, dipaparkan kesimpulan utama yang secara langsung menjawab rumusan masalah mengenai perancangan dan efektivitas sistem asesmen kepribadian hibrida. Selanjutnya, bab ini juga memberikan saran-saran konstruktif yang dapat dijadikan dasar untuk eksplorasi penelitian dan pengembangan sistem lebih lanjut di masa mendatang.
% Ubah bagian-bagian berikut dengan isi dari penutup

\section{Kesimpulan}
\label{sec:kesimpulan}

Berdasarkan hasil eksperimen penelitian dan analisis yang dilakukan, berikut adalah kesimpulan yang didapatkan:

\begin{enumerate}[nosep]

  \item Implementasi model regresi berbasis arsitektur Transformer untuk memprediksi skor kepribadian Big Five dilakukan melalui modifikasi komponen \textit{custom regression head}. Komponen ini menggunakan kombinasi lapisan linear berdimensi 512 neuron, fungsi aktivasi \textit{non linear} GELU, \textit{dropout} sebesar 30\%, dan fungsi sigmoid untuk memetakan ekstraksi fitur visual menjadi rentang probabilitas regresi yang valid. Berdasarkan evaluasi terhadap enam varian model, konfigurasi ViT-B/16 AugReg (IN21k) yang dioptimalkan dengan skenario \textit{arch tuning} beserta augmentasi menjadi model dengan kinerja terbaik. Model ini tidak hanya mampu mencapai tingkat presisi tertinggi dengan rata-rata akurasi puncak sebesar 0,911, tetapi juga mencapai nilai \textit{coefficient of determination} ($R^2$) sebesar 0,416 pada data uji independen. Pencapaian metrik $R^2$ ini menunjukkan bahwa penambahan kompleksitas arsitektur dan variabilitas data berhasil mencegah \textit{overfitting} sekaligus memungkinkan model menangkap keunikan representasi kepribadian individu.

  \item Integrasi hasil prediksi regresi ke dalam LLM untuk menghasilkan interpretasi deskriptif dilakukan melalui sistem \textit{prompt engineering} multimodal. \textit{Prompt} dirancang menggunakan empat blok informasi utama. \textit{System instruction} untuk menentukan batasan persona, \textit{visual context} berupa citra dengan \textit{encoding} Base64, \textit{numerical context} (skor kepribadian Big Five), dan \textit{task instruction}. Pengujian memberikan hasil bahwa penerapan gaya \textit{roleplay prompting} sangat efektif dalam mentransformasi narasi LLM khususnya pada keluarga model Gemma 4 yang kaku dan analitis menjadi teks profil psikologis yang jauh lebih hangat, empatik, dan interaktif layaknya sebuah sesi konseling. Meskipun demikian, evaluasi \textit{LLM-as-a-judge} menunjukkan bahwa instruksi \textit{standard prompting} cenderung menghasilkan struktur laporan klinis yang lebih terstruktur poin per poin, yang membuat penilaian metrik evaluator G-Eval lebih tinggi dibandingkan \textit{roleplay prompting}.

  \item Kinerja sistem asesmen kepribadian hibrida secara keseluruhan menunjukkan kemampuan analisis yang menyeluruh dari tahap awal hingga akhir terlepas dari adanya keterbatasan pada modalitas visual. Pada tahap prediksi visual menggunakan model prediksi Transformer, sistem menunjukkan kinerja tertinggi dalam mengenali sifat \textit{conscientiousness} dan \textit{extraversion} yang terlihat kuat pada fitur fisik, namun mengalami tantangan pada sifat \textit{agreeableness} yang bergantung pada interaksi sosial. Selanjutnya pada tahap interpretasi akhir (generasi teks), hasil G-Eval menunjukkan tidak adanya LLM yang mendominasi seluruh metrik penilaian. Oleh karena itu, sistem ini mengimplementasikan pendekatan dua model alternatif dengan menetapkan Qwen3-VL-32B-Instruct sebagai model utama karena meraih rata-rata keseluruhan tertinggi (0,9136) dengan narasi persona psikolog yang kaya, dan Gemma-4-31B-IT sebagai alternatif karena menawarkan keseimbangan performa yang lebih stabil antara akurasi analitis (0,9504) dan koherensi visual, serta kepatuhan yang tinggi terhadap batasan instruksi format. Untuk memvalidasi kualitas narasi interpretasi tersebut secara lebih mendalam, dilakukan pula evaluasi langsung oleh pakar psikolog terhadap kedua model dan gaya respons yang dipilih. Hasil evaluasi ini menunjukkan skor rata-rata yang jauh lebih rendah, yaitu pada kisaran 52\% hingga 61\% dari skor maksimal pada empat aspek penilaian, dibandingkan skor G-Eval yang secara konsisten berada di atas 79\% pada kombinasi yang sama. Hasil ini mengindikasikan bahwa meskipun G-Eval memberikan skor tinggi terhadap narasi interpretasi yang dihasilkan dari segi struktur maupun kelengkapan jawaban, kualitas tersebut belum memenuhi standar interpretasi psikolog profesional dari segi ketepatan klinis dan keterhubungan dengan fitur visual pada gambar. Ditemukan pula bahwa pengaruh gaya \textit{roleplay prompting} terhadap kualitas interpretasi tidak seragam antar model. Skor mengalami peningkatan pada Gemma-4-31B-IT namun menurunkan skor pada Qwen3-VL-32B-Instruct, sehingga kinerja sistem secara keseluruhan turut bergantung pada kesesuaian gaya \textit{prompting} dengan karakteristik masing-masing model.

\end{enumerate}

\section{Saran}
\label{chap:saran}

Berdasarkan penelitian yang dilakukan, terdapat kekurangan atau potensi yang dapat dieksplorasi lebih lanjut untuk pengembangan selanjutnya antara lain:

\begin{enumerate}[nosep]

  \item Penelitian selanjutnya disarankan untuk mengadopsi modalitas data lain secara penuh, seperti integrasi analisis gerak spasio-temporal dari video dan analisis fitur suara. Penambahan modalitas ini diharapkan mampu mengatasi keterbatasan penggunaan gambar statis, khususnya dalam memprediksi sifat kepribadian yang sangat berkaitan dengan cara manusia berinteraksi dan berempati seperti \textit{agreeableness}.

  \item Diperlukan eksplorasi terkait \textit{parameter tuning} tingkat lanjut pada fase inferensi LLM seperti penyesuaian \textit{temperature}, \textit{top-p}, atau \textit{frequency penalty} untuk mengontrol panjang respons teks dengan lebih presisi. Hal ini diperlukan untuk mengatasi perilaku beberapa LLM yang terlalu elaboratif seperti Qwen3-VL sehingga rentan melanggar instruksi batasan kata dan menyebabkan hasil interpretasi terpotong secara paksa karena menyentuh batas maksimum token sistem.

  \item Evaluasi oleh psikolog pada penelitian ini masih terbatas pada satu orang penilai dengan 20 sampel narasi. Penelitian selanjutnya disarankan untuk melibatkan lebih dari satu psikolog penilai dengan cakupan sampel yang lebih luas, guna memperoleh hasil validasi yang lebih representatif dan dapat diandalkan secara statistik.
  
  \item Kriteria dan instrumen penilaian kualitas interpretasi kepribadian pada penelitian ini disusun berdasarkan kebutuhan praktis penelitian dan belum disesuaikan dengan teori psikometri yang telah teruji secara ilmiah, sebagaimana turut disampaikan oleh psikolog penilai pada proses evaluasi. Penelitian selanjutnya disarankan untuk menyusun instrumen penilaian yang lebih baku dengan menggunakan ilmu psikometri, termasuk memperhatikan prinsip validitas dan reliabilitas dalam proses penyusunannya, sehingga evaluasi kualitas interpretasi kepribadian yang dihasilkan lebih terukur dan memiliki dasar ilmiah yang lebih kuat.

\end{enumerate}
